\begin{table}[H]
\small
\setlength{\tabcolsep}{5pt}
\renewcommand{\arraystretch}{0.95}

\caption{\label{tab:country-heterogeneity-cumulative-gap}Cumulative inflation gap between the first and fifth income quintiles, 2020--2024}
\centering
\begin{tabular}[t]{llll}
\toprule
Country & Q1 cumulative inflation & Q5 cumulative inflation & Q1-Q5 gap\\
\midrule
Austria & 23.4 & 24.0 & -0.7\\
Cyprus & 17.8 & 17.5 & 0.2\\
Estonia & 46.6 & 39.8 & 6.9\\
Finland & 13.6 & 15.9 & -2.2\\
France & 17.8 & 17.9 & -0.1\\
\addlinespace
Germany & 20.0 & 22.5 & -2.5\\
Greece & 17.9 & 17.6 & 0.3\\
Ireland & 21.4 & 16.0 & 5.5\\
Italy & 19.8 & 19.8 & -0.0\\
Latvia & 38.1 & 32.1 & 6.1\\
\addlinespace
Lithuania & 38.2 & 35.5 & 2.7\\
Luxembourg & 16.6 & 17.8 & -1.1\\
Netherlands & 22.3 & 23.6 & -1.3\\
Portugal & 17.9 & 18.0 & -0.1\\
Slovakia & 31.2 & 32.3 & -1.1\\
\addlinespace
Spain & 18.3 & 18.5 & -0.2\\
\midrule
\textbf{Euro area} & \textbf{19.8} & \textbf{20.6} & \textbf{-0.7}\\
\bottomrule
\end{tabular}
\vspace{-0.4em}
\begin{flushleft}
\footnotesize{\emph{Notes:} Entries are percentage points. Cumulative inflation is computed from average income-quintile price indices in 2020 and 2024, using RAS expenditure weights. France is computed at COICOP level 3 with INSEE income quintiles; the euro-area row is rebuilt from country-level indices and EA20 HICP country weights using this France level-3 replacement. Other country rows use the level-2 income-quintile setup. The gap is Q1 minus Q5. \emph{Sources:} Eurostat; authors' calculations.}
\end{flushleft}
\end{table}
